\documentclass[10pt,a4paper]{amsart}
\usepackage[margin=19mm]{geometry}
\usepackage{amsmath,amssymb,mathtools}
\usepackage[scr=rsfs,cal=euler]{mathalfa}
\usepackage{xfrac}
\usepackage[unicode,hidelinks,bookmarks=false]{hyperref}
\newcommand{\ie}{i.e.\ }
\newcommand{\cf}{cf.\ }
\newcommand{\resp}{respectively\ }
\newcommand{\Subsection}{\subsection}
\providecommand{\nobreakdash}{\nobreak\hbox{-}}
\setlength{\emergencystretch}{2em}
\multlinegap0pt
\usepackage{xcolor}
\usepackage[normalem]{ulem}
\usepackage{tikz}
\usepackage{amssymb}
\usepackage{mathtools}
\newtheorem{thm}{Theorem}
\newtheorem{prop}{Proposition}
\newcommand{\Aut}{{\operatorname{Aut}}}
\newcommand{\F}{\mathbb{F}}
\newcommand{\ord}{{\operatorname{ord}}}
\title{The Lang-Trotter conjecture on average for genus-$2$ curves with $S_3$ reduced automorphism group}
\author{Chihiro Ando, Shushi Harashita}
\date{Source: arXiv:2604.00822v1 (CC BY 4.0)}
\begin{document}
\maketitle

\noindent\emph{Source and licence.} The Lang-Trotter conjecture on average for genus-$2$ curves with $S_3$ reduced automorphism group. Version arXiv:2604.00822v1 (CC BY 4.0), distributed by arXiv under the Creative Commons Attribution 4.0 International licence, http://creativecommons.org/licenses/by/4.0/, as recorded in the arXiv licence metadata for this exact version and shown on the abstract page https://arxiv.org/abs/2604.00822v1. Full licence notice: \texttt{licenses/arxiv-2604.00822-CC-BY-4.0.txt}.\par\smallskip

\noindent\textit{Editorial scope.} Counted unit: the article's prop GZmultiplicity (source0.tex lines 454--494). The article's complete original proof of it is delivered with it (source0.tex lines 495--538, one \texttt{proof} environment of 44 lines). No mathematics is written, repaired or re-derived: the statement and the proof are the article's own lines, in the article's own notation, and no retained line is substituted or re-flowed.  Retained uncounted as the source-local context the proof needs: lines 353--372; lines 357--367.  Nothing was omitted from any retained span, so the package carries no omission record.  No retained line contains a citation, so the wrapper carries no bibliography and no bibliography entry.  No retained line contains a figure environment, an image-inclusion command or drawing code, so no image is delivered and the package declares no asset dependency.  Editorial formatting only, in addition: an AMS \texttt{amsart} preamble that restates the article's own declarations and macros used by the retained lines, namely the environments \texttt{thm}, \texttt{prop} on separate counters, together with the standard packages those lines need.  The retained environments print in this document as Theorem 1, Proposition 1 in order of appearance, whereas the article prints the same items with the numbering of its own full document.  The complete native source of the article as distributed in the arXiv e-print for this exact version is bundled unchanged as \texttt{sources/197612/source0.tex}.
\begin{thm}[\textcolor{red}{To be checked}]\label{thm:GZ_for_3p}
Assume  that $-D$ is fundamental
and that $p \equiv 1 \pmod 4$ with $(D,3p) = 1$ and $\left(\frac{-D}{p}\right)\ne 1$. 
Then  
\begin{equation}\label{eq:GZ_for_3p}
\frac{1}{w}
\sum_{j}
%_{
%\begin{aligned}
%    [E]\text{\rm : CM of type }D\\
% j(E_p) \in \F_p
%\end{aligned}
% }
m_j\cdot \#\Aut(E_j)=    \ord_p J(-D, -3p)^2
\end{equation}
with the summation over distinct roots $j\in\F_{p^2}$ of the greatest common divisor of $P_D(X) \bmod p$ and $P_{3p}(X) \bmod p$, where $m_j$ is the product of the multiplicities at $j$ of $P_D(X) \bmod p$ and $P_{3p}(X) \bmod p$.
%the reduction of $E$ modulo $p$ and
%$E$ runs over the isomorphism classes of
%elliptic curves of CM type $D$ with $j(E_p)\in\F_p$.
\end{thm}

\begin{equation}
\frac{1}{w}
\sum_{j}
%_{
%\begin{aligned}
%    [E]\text{\rm : CM of type }D\\
% j(E_p) \in \F_p
%\end{aligned}
% }
m_j\cdot \#\Aut(E_j)=    \ord_p J(-D, -3p)^2
\end{equation}

\begin{prop}
\label{GZmultiplicity}
\begin{enumerate}
\item[(1)] For $D=8$ and $j=8000$ (the root of $P_D(j)$),
we have
\[
m_j = \begin{cases}
4 & \text{ if } p \equiv 5 \pmod 8 \text{ and } p \ne 5,\\
2 & \text{ if } p = 5,\\
0 & \text{ otherwise. }\\
\end{cases}
\]
%$m_j = 2$ for $p = 5$ and $m_j = 4$ for $p > 5$ with $p\equiv 5\pmod 8$. Otherwise $m_j=0$.
\item[(2)] For $D=20$ and a root $j$ of $P_D(j)$, we have
\[
m_j = \begin{cases}
2 & \text{ if } p \equiv 0, 2,3 \pmod 5 \text{ and } p\ne 13,\\
%2 & \text{ if } p=5\\
4 & \text{ if } p=13,\\
0 & \text{ otherwise.} \\
\end{cases}
\]
%\item[(3)] (non fundamental, extra arguments) For $D=32$ 
%and a root $j$ of $P_D(j)$, we have
%\[
%m_j = \begin{cases}
%2 & \text{ if } p \equiv 5 \pmod 8 \text{ and } p\ne 13,\\
%4 & \text{ if } p = 13,\\
%0 & \text{ otherwise. }\\
%\end{cases}
%\]
\item[(3)] For $D=35$, we have
\[
m_j = \begin{cases}
2 & \text{ if } p = 5, 37, 41, 53, 89, 101,\\
4 & \text{ if } p = 61,\\
0 & \text{ otherwise. }\\
\end{cases}
\]
\end{enumerate}
\end{prop}

\begin{proof}
(1) First consider the case that $p>5$.
We have $F(m) \ne 0$ with $m=\frac{24p-x^2}{4}$ only if $x=0$
and $p\equiv 5 \pmod 8$. Furthermore $F(6p) = 4$ by $\varepsilon(p)=-1$, $\varepsilon(2)=1$ and $\varepsilon(3)=1$. 
For $j=8000$, Theorem \ref{thm:GZ_for_3p} reads
\[
\frac{1}{2} m_j \# \Aut(E_j) = 4.
\]
Since $\# \Aut(E_j) = 2$,
we have $m_j = 4$ for $p\equiv 5 \bmod 8$ with $p > 5$.

Next consider the case that $p=5$.
We have $F(m)\ne 0$ with $m=\frac{24p-x^2}{4}$ if $x=0$, $\pm 2p$.
Furthermore $F(6p) = 4$ and $F(p) = 1$. Hence
Theorem \ref{thm:GZ_for_3p} reads
\[
\frac{1}{2} m_j \# \Aut(E_j) = 4+1+1.
\]
Since $\# \Aut(E_j) = 6$,
we have $m_j = 2$ for $p=5$.

(2) 
First consider the case of $p>17$.
The right hand side of \eqref{eq:GZ_for_3p} for $D=20$ is $\ord_p F(15p)$,
as the contribution comes only from $x=0$.
If $p \equiv 1,4 \pmod 5$, then $\varepsilon(5)=\left(\frac{-3p}{5}\right)=-1$,
whence $\ord_p F(15p)=0$. Thus we assume $p \equiv 2,3 \pmod 5$.
Note
\[
P_{20}(X)=X^2 - 1264000X - 681472000.
\]
The discriminant of $P_{20}(X)$ is $2^{18} 5^3 13^2 17^2$.
By the law of quadratic reciprocity, $\displaystyle \left(\frac{5}{p}\right) = \left(\frac{p}{5}\right)$, which is $-1$ by $p \equiv 2,3 \pmod 5$. Hence
$P_{20}(X)$ considered as a polynomial over ${\mathbb F}_p$ is irreducible. %since $p \equiv 2,3 \pmod 5$.
Hence the multiplicity $m_j$ of a root $j$ of $P_{20}(X)$ is independent of the choice of the root. Theorem \ref{thm:GZ_for_3p} reads
\[
\frac{1}{2} \sum_{j \text{ s.t. } P_{20}(j)=0} m_j \# \Aut(E_j) = \ord_p F(15p) = 4,
\]
since $\varepsilon(3)=\varepsilon(5)=1$. By $\# \Aut(E_j)=2$,
we have $m_j = 2$. The number of the remaining primes (i.e., $p\le 17$) are finite. The lemma for the cases can be confirmed through exhaustive calculations.

(3) Put $x=py$. Then $3pD-x^2 = p(105-py^2)>0$. The set of pairs $(p,x)$ is finite. The lemma
can be confirmed through exhaustive calculations.
\end{proof}

\end{document}
